\chapter{Conclusiones provisionales}

\begin{enumerate}
\item Se implementó una arquitectura reproducible en ROS~2 para generar y
supervisar postura, marcha paso y gateo del Nova Spot Micro en Gazebo y MuJoCo.
El resultado aporta evidencia parcial a los tres primeros objetivos, pero no
cierra la caracterización ni la integración física.

\item La corrección del planificador temporal alineó la cadencia ejecutada con
los 4,32~s configurados por ciclo. Dos ensayos de la nueva versión presentaron
diferencias menores al 0,2~\% en las métricas medias seleccionadas, lo que
respaldó una línea base repetible en simulación.

\item La marcha paso se ejecutó con pose, IMU, contactos y articulaciones bajo
un contrato común en Gazebo y MuJoCo. Once ciclos mostraron diferencias de
avance y contacto: se cerró la equivalencia instrumental, no la dinámica.

\item El ensayo de contacto crudo y filtrado mostró que las interrupciones de
RL y RR tuvieron máximos inferiores a 0,09~s y nunca superaron la persistencia
de 0,12~s. Por tanto, no se demostró un vuelo trasero sostenido. El avance del
robot no fue evidencia suficiente del cumplimiento de la secuencia ideal de
contactos.

\item El modelo actual permitió calcular referencias, pose de pies y margen de
estabilidad, pero continuó siendo nominal y no identificado. No corresponde
denominarlo gemelo digital identificado ni extrapolar directamente sus métricas
al hardware; se mantiene la denominación modelo digital configurable.

\item No existe todavía evidencia para concluir que una política de aprendizaje
por refuerzo mejore la estabilidad o repetibilidad del robot. Los objetivos
cuarto y quinto permanecen abiertos hasta entrenar las políticas, validar la
marcha física y ejecutar la comparación emparejada prevista.
\end{enumerate}

Estas conclusiones son provisionales y se actualizarán únicamente después de
incorporar resultados físicos y la comparación nominal frente a nominal más
aprendizaje por refuerzo.
